\documentclass{article}
\usepackage[UTF8]{ctex}
\usepackage{amsmath,amssymb,amsthm,bm}
\usepackage[a4paper,margin=1in]{geometry}

\title{模型说明与可微替代损失}
\author{}
\date{}

\begin{document}
\maketitle

\section{问题定义}

设输入特征为 \(x\)，模型参数为 \(\theta\)。模型首先根据输入特征输出预测值
\begin{equation}
\hat y = f_\theta(x).
\end{equation}

在得到预测值 \(\hat y\) 后，再通过一个决策函数 \(h(\cdot)\) 生成最终决策
\begin{equation}
\hat z = h(\hat y).
\end{equation}

其中，\(\hat z\) 表示由模型预测结果诱导出的决策。通常，\(h(\cdot)\) 是一个基于打分进行最优选择的决策规则，例如通过 \(\arg\max\) 从若干候选动作中选出收益最大的动作。

\section{原始决策损失}

我们希望模型训练目标能够直接服务于最终决策效果，而不是仅仅提高预测精度。为此，定义基于决策结果的损失函数。记真实环境下与决策相关的指标为 \(y\)，随机分配到的决策标签为 \(t\)。

原始损失可以写为
\begin{equation}
\mathcal{L}_{\text{orig}}(\theta) = - \mathrm{EOM}\bigl(h(\hat y)\bigr).
\end{equation}

进一步地，若将 \(\mathrm{EOM}\) 写成对满足决策匹配条件样本的求和形式，则可表示为
\begin{equation}
\mathcal{L}_{\text{orig}}(\theta)
= - \sum_{i} y_i \, \mathbf{1}\!\left(h(\hat y_i)=t_i\right),
\end{equation}
其中：
\begin{itemize}
    \item \(i\) 表示样本索引；
    \item \(y_i\) 表示样本 \(i\) 的真实结果；
    \item \(t_i\) 表示样本 \(i\) 被随机分配到的决策标签；
    \item \(\mathbf{1}(\cdot)\) 为示性函数，当括号内条件成立时取 1，否则取 0。
\end{itemize}

该损失的含义是：只有当模型产生的决策 \(h(\hat y_i)\) 与样本实际被分配到的决策标签 \(t_i\) 一致时，样本 \(i\) 的观测结果 \(y_i\) 才会被计入目标函数。

\section{不可微问题}

在实际中，决策函数 \(h(\hat y)\) 常写为
\begin{equation}
h(\hat y)=\arg\max_{k \in \mathcal{A}} s_k(\hat y),
\end{equation}
其中 \(\mathcal{A}\) 表示候选动作集合，\(s_k(\hat y)\) 表示动作 \(k\) 在预测值 \(\hat y\) 下对应的打分函数。

由于 \(\arg\max\) 是离散选择操作，因此 \(h(\hat y)\) 对 \(\hat y\) 不可微，从而损失函数
\(\mathcal{L}_{\text{orig}}(\theta)\) 无法直接通过梯度下降对参数 \(\theta\) 进行优化。

\section{Softmax 可微替代}

为了实现端到端训练，这里采用最简单的可微近似方式：用 softmax 替代原始的 \(\arg\max\) 操作。

首先，对每个候选动作 \(k\) 定义其得分
\begin{equation}
s_k = s_k(\hat y).
\end{equation}

然后构造 softmax 概率：
\begin{equation}
\pi_k(\hat y)
=
\frac{\exp(s_k/\tau)}
{\sum_{j \in \mathcal{A}} \exp(s_j/\tau)},
\end{equation}
其中 \(\tau>0\) 为温度参数。当 \(\tau\) 越小时，softmax 分布越接近 one-hot，从而越接近原始的 \(\arg\max\) 选择。

于是，原来的硬决策
\begin{equation}
h(\hat y)=\arg\max_{k \in \mathcal{A}} s_k(\hat y)
\end{equation}
可以用 softmax 概率向量
\begin{equation}
\tilde z(\hat y)=\bigl(\pi_k(\hat y)\bigr)_{k\in\mathcal{A}}
\end{equation}
进行替代。

\section{替代损失函数}

将原先基于硬匹配的示性函数
\(\mathbf{1}(h(\hat y_i)=t_i)\)
替换为 softmax 下对应动作 \(t_i\) 的概率 \(\pi_{t_i}(\hat y_i)\)，可得到可微替代损失：
\begin{equation}
\mathcal{L}_{\text{soft}}(\theta)
=
- \sum_i y_i \, \pi_{t_i}(\hat y_i).
\end{equation}

将 \(\hat y_i=f_\theta(x_i)\) 代入，上式也可写为
\begin{equation}
\mathcal{L}_{\text{soft}}(\theta)
=
- \sum_i y_i \,
\frac{
\exp\!\bigl(s_{t_i}(f_\theta(x_i))/\tau\bigr)
}{
\sum_{j\in\mathcal{A}}
\exp\!\bigl(s_j(f_\theta(x_i))/\tau\bigr)
}.
\end{equation}

\section{训练解释}

该替代损失具有以下特点：

\begin{itemize}
    \item 当模型更倾向于选择与随机分配标签 \(t_i\) 一致的动作时，\(\pi_{t_i}(\hat y_i)\) 会增大；
    \item 若对应样本的真实结果 \(y_i\) 较大，则该样本会对目标函数产生更强的正向贡献；
    \item 整个目标函数关于 \(\hat y\) 以及参数 \(\theta\) 都是可微的，因此可以利用反向传播进行优化。
\end{itemize}

\section{总结}

综上，该模型的流程可以概括为：
\begin{equation}
x \xrightarrow{f_\theta} \hat y \xrightarrow{h} \hat z.
\end{equation}

原始损失基于硬决策选择：
\begin{equation}
\mathcal{L}_{\text{orig}}(\theta)
=
- \sum_i y_i \, \mathbf{1}\!\left(h(\hat y_i)=t_i\right),
\end{equation}
但由于 \(h(\hat y)\) 中含有 \(\arg\max\)，损失不可微。

因此，采用 softmax 对 \(\arg\max\) 进行平滑近似，得到可微替代损失：
\begin{equation}
\mathcal{L}_{\text{soft}}(\theta)
=
- \sum_i y_i \, \pi_{t_i}(\hat y_i),
\end{equation}
从而可以将预测模型与决策过程联合起来进行端到端训练。

\end{document}